\section{Paso 10: conclusiones y recomendaciones}
\label{sec:paso10-conclusiones}
% Origen: regla de decisión de la sección 8 (fijada antes de correr), resultados de la sección 10 y limitaciones
% documentadas en 02-modelo-anylogic.md.

\subsection{Respuesta a la pregunta del estudio}

\pendiente{responder con los intervalos de la tabla~\ref{tab:diferencias}: si habilitar los 22 molinetes
(\EI{}), validar más rápido (\ET{}) o las dos cosas (\EIT{}) reducen la espera de la hora pico frente a la
situación actual, y cuánto.}

\subsection{Recomendación}

\pendiente{aplicar la regla de decisión: recomendar la alternativa de menor P90 cuyo intervalo frente a las otras
lo sostenga, o describir el compromiso si las primarias no coinciden.}

\subsection{Alcance y limitaciones}

Las conclusiones valen dentro de los supuestos del modelo. Los que más pueden pesar son estos:
\begin{itemize}
  \item \textbf{Espera no validada.} No hay observaciones de la espera ni del largo de cola en Constitución
        (sección~\ref{sec:paso6-validacion}). Los niveles absolutos dependen de los supuestos; la comparación
        entre escenarios es más robusta, porque todos los comparten.
  \item \textbf{Reparto más parejo que el real.} El modelo reparte la carga entre molinetes más parejo que la
        gente real. Si los molinetes que agrega \EI{} quedaran en una posición poco atractiva, se usarían menos y
        el modelo sobreestimaría su efecto.
  \item \textbf{Proporción del Roca.} Que el 80~\% de las validaciones del vestíbulo Principal venga del Roca es
        un supuesto: los datos disponibles no permiten estimarla.
  \item \textbf{Tiempo de validación rápida.} \ET{} y \EIT{} usan el extremo optimista del rango técnico, no una
        medición del molinete EMV o QR.
  \item \textbf{Geometría del subsuelo.} La plaza de transbordo, el ancho del pasillo y el interior del
        vestíbulo son supuestos; el hall, los andenes y las escaleras mecánicas siguen el plano oficial. Las
        escaleras fijas del bloque central no se modelan.
  \item \textbf{Arranque.} El sistema empieza vacío a las 07:00 y no genera el pasaje de los trenes anteriores,
        así que la primera ventana queda con menos gente que la real.
  \item \textbf{Pasajeros fuera del modelo.} Solo se modela a quien entra al vestíbulo Principal. Los que bajan
        del Roca y salen a la calle o van a otro acceso no se simulan, aunque compartan el hall.
\end{itemize}

\subsection{Trabajo futuro}

\begin{itemize}
  \item Medir en el lugar la espera y el largo de cola en la hora pico, para validar los niveles absolutos.
  \item Incorporar el plano del subsuelo de SBASE, si se recibe, y recalibrar la elección de molinete con el
        reparto observado.
  \item Medir el tiempo de validación del molinete EMV o QR que ya existe en la estación.
  \item Evaluar E2, el desvío de parte del flujo al vestíbulo Plaza, que exige modelar ese segundo circuito.
  \item Agregar un período de calentamiento desde las 06:45 para representar la primera ventana.
\end{itemize}
